\section{File System Object Isolation}\label{sec:ref-mon}\label{sec:design:ref-mon}\label{sec:design:file-path}

% \begin{figure}[t]
%     \centering
%     \includegraphics[width=.98\columnwidth]{images/capacity-syscall-filter}
%     \caption{\thename syscall filter vs. other syscall filter designs when filtering \texttt{sys\_read} and \texttt{sys\_write}.}
%     \label{fig:ref-mon}
% \end{figure}
% \chcomment{Rule-based filtering vs. capability-based filtering}


% Also, importing and using file resources must be mediated to prevent leaking % sensitive assets of a protected domain~\cite{jenny}. \thename's reference monitor enables the complete mediation of file path and file descriptor usages by authenticating the \gls{syscall} arguments. 
% (e.g., \cpath into \fd)

% Previous works on \gls{ipi} implemented reference monitors that mediate syscalls, along with their memory isolation scheme~\cite{erim,hodor,jenny,cerberus,pku-pitfall}. The necessity of such monitors arises because the kernel is unaware of the in-process domains, and adversaries in the untrusted domain may launch various forms of confused deputy attacks on the trusted domain. 

% \thename references also serve as a non-forgeable token in a capability-based access control scheme.

\keypoint{
	\thename includes a reference monitor that verifies the \emph{arguments} of syscalls to enforce the complete mediation (\srcompmed) of system resource accesses.
	This capability-inspired authentication mechanism drastically differs from previous works implementing \gls{acl}-based reference monitors for in-process syscall filtering and file access control~\cite{erim,hodor,jenny,cerberus,pku-pitfall}.
}
In those systems, enforcing access control on domain-private sensitive file objects requires a separate logic for the monitor detached from the program semantics.
Moreover, \gls{acl}-based access control requires keeping track of the isolated objects and the domains that can access them at runtime~\cite{jenny}.
In contrast, the \fds in \thename are cryptographically secured; the proof of ownership and object attributes are attached to the reference itself, regardless of the number of subjects sharing a resource.

We implement the reference monitor in about $1000$ LoC. It provides domain key switching and \gls{syscall} interception through \gls{syscall} table hooking using a similar approach to a previous efficient reference monitor~\cite{jenny}.
That is, we employ a kernel module that hooks the \gls{syscall} table with custom authentication/signing routines that are invoked before and after a \gls{syscall}'s execution.
% as a kernel module that hooks the \gls{syscall} table and prepend domain authentication routines using the \gls{pa} instructions. This method allows our reference monitor with minimal overheads. Note that the method is similar t
To avoid the system monitor affecting all processes within the system, we set an unused flag in \texttt{struct thread\_info->flags} to mark \thename-enabled processes during initialization and check for the flag before performing syscall authentication.


% Note that we do not seek completeness with the syscall filtering policy in our reference monitor.
% The syscall APIs were never developed with intra-process isolation in mind. This large attack surface entails various domain-to-domain \emph{confused deputy} problems that undermine the security of protected domains.
% Many such attacks have been discussed in the previous works~\cite{pku-pitfall,jenny,cerberus}.
% However, we regard the issue as an open challenge and  leave the completeness of the syscall filtering out of scope.

% \keypoint{
% 	\thename supports the authentication of file system paths (\cpath) and file descriptors (\fd).
% }



% As demonstrated later, we utilize the attributes to attach Capsicum-inspired policies that limit the capabilities of file descriptors to be used in syscalls.
%
% Also, with hardware acceleration, \thename's system object authentication only requires minimal kernel extensions and can be applied at the intra-process granularity.
%

% is performed automatically with \thename's facilities.  
% \thename's cryptographically-protected references are intuitively built into the program as values, 


% \noh{
% Our capability-based resource filtering offers several advantages.
% First, capabilities allow users to control authority per resource, providing a more granular access control level.
% Although it may be possible to achieve the same level of granularity with a traditional access control system, capabilities reduce complexity by only requiring a corresponding authority token, regardless of the number of subjects sharing a particular resource.
% In contrast to a syscall filtering, which limits the role or privilege of a specific domain, our approach uses resource filtering to constrain what can be accessed rather than what can be performed.
% As a result, our capability-based resource filtering offers a more flexible and secure approach to access control.
% }



% By adopting capability-based access control, \thename enables drastically simplified authentication and management of sensitive file object references in syscalls.
% We demonstrate the adv    antages of \thename's capability-based reference monitor design over \gls{acl}-based reference monitors. 


% \hdr{In-program and coherent sensitive object access control} 
% Another difference is the programming model between the two monitor designs.
% Enforcing ACL on sensitive objects requires a separate code for the monitor, which is detached from the program semantics.
% In contrast, cryptographically-secured object references in \thename are coherently built into the program through a \thename library, in-kernel monitor, and compiler instrumentation.
% Managing the access control on sensitive objects in \thename is achieved through in-place annotations.



% \added{
% In \thename, file object isolation is provided by the complete mediation (\srcompmed) of resource usage provided by the reference monitor.
% Isolation of \cpath references is provided by resolving the path to get its underlying file object (\ie, \texttt{struct inode}), then authenticate its attached the \gls{acl}-like data structure, as explained in (\cref{sec:design:path}).
% File descriptors isolation is provided with Capsicum-inspired capability-enhanced file descriptors~\cite{capsicum} that are efficiently constructed and verified with \gls{pa} hardware acceleration (\cref{sec:design:fd}).
% The reference monitor also propagates ownership of references when they change forms, which enables the life-cycle isolation of program resources.
% }

% The reference monitor of \thename performs several tasks: First, it enables the authenticated domain switching mechanism described in \cref{sec:design:impers} by serving the \texttt{capac\_enter} and \texttt{capac\_exit} requests from userspace.
% Second, it also enforces the complete mediation (\srcompmed) of the file path and file descriptor uses by interposing \fd/\cpath-accepting syscalls and authenticating those references inside arguments with \gls{pa}.



% \added{
% \autoref{fig:ref-mon} presents a qualitative comparison between \thename and other technologies that filter syscalls.
% Specifically, we compare \thename with two existing designs: the industry-standard seccomp-based filter and the user-delegate technique proposed recently~\cite{jenny}.
% The seccomp-based filter relies on a set of predefined rules implemented in the kernel to restrict a program's access to syscalls. To use seccomp, a program must define a filter that specifies the filtering logic, typically using the BPF language.
% The filter is then loaded into the kernel, which applies it to the program's syscalls to determine whether they are authorized.
% However, seccomp-based mechanisms have high overheads and are inflexible, making them less ideal for intra-process isolation scenarios~\cite{jenny}.
% To address these limitations, the user-delegate-based syscall filter was proposed~\cite{jenny}.
% The design delegates the task of authorizing syscalls to an in-process reference monitor, allowing for more flexible filtering logic to be registered for nested domains while reducing syscall authorization overhead.
% }

% \added{
% \thename's capability principle simplifies enforcement by authenticating references with the current domain authentication key.
% Similar to the user-delegate-based filter, our approach requires minimal kernel extension, as the reference contains all the necessary information to authorize resource usage.
% However, while each domain in user-delegate requires a dedicated filter, all \thename domains use a uniform authentication logic, requiring only a simple \gls{pa} instruction invocation on the authenticated reference.
% Furthermore, \thename employs the same authentication logic to filter both read and write syscalls since they both accept a file descriptor as an argument.
% This is not the case with seccomp-based or user-delegate-based filters. 
% }



% For instance,...

% as a kernel module that hooks the \gls{syscall} table and prepend domain authentication routines using the \gls{pa} instructions. This method allows our reference monitor with minimal overheads. Note that the method is similar t


% \hdr{Programming interfaces}

% \added{
% A user employs \thename to isolate file system resources through \libcapac APIs, namely \texttt{capac\_init}, \texttt{capac\_delegate\_fd} and \texttt{capac\_limit\_fd}, which are described in \autoref{tab:apis}.
% Notably, the protection of file system object is kickstarted by the \texttt{capac\_init} function call.
% The function receives a configuration structure, including a file path-owner domains mapping and other necessary information.
% \texttt{capac\_init} then transitions the program into an irreversible \emph{protected mode}, where the reference monitor authenticates every \cpath and \fd arguments used in \glspl{syscall}.
% Afterward, \thename protection coherently exists in the program, with file descriptors and path being transparently turned into unforgeable capability tokens used to request access to resources. 
% }

% \texttt{capac\_init} performs the initialization of \thename facilities, such as initializing the kernel-managed domain keys. 
% It also receives from the user a list of domain-private \cpath and their owner domains, which it uses to establish the file-owner relationship (described later in \cref{sec:design:path}). 

% \added{
% The \texttt{capac\_delegate\_fd} API provides a \emph{runtime} mechanism for domains to share access to domain-private file descriptors.
% When a delegation request is received, the reference monitor authenticates the delegated \fd to verify that its owner is requesting delegation.
% It then re-signs the reference with the target's domain key and the requested context, then returns the newly signed \fd to the userspace.
% }

% \added{
% Finally, \texttt{capac\_limit\_fd}
% }


% \subsection{Authenticating file system paths}\label{sec:design:path}
\subsection{Enforcing domain-private file-system paths}
\label{sec:design:path}
% \issue{Maybe remove figure 5}
% \begin{figure}[b]
% 	\centering \includegraphics[width=0.90\columnwidth]{images/capacity-path-filter}
% 	\caption{\thename file path authentication logic}
% 	\label{fig:path}
% \end{figure}
\keypoint{
	\thename supports file path isolation to be compatible with existing programs without significant rewriting effort.
}
\thename's file path authentication utilizes the pre-existing methods for attaching policies to file system objects used by kernel security subsystems (e.g., SELinux~\cite{selinux}).

File paths protection is kickstarted by the \texttt{capac\_init} API call that sends a list of domain-private files and their owners to the reference monitor.
The reference monitor first resolves each path to obtain the corresponding \texttt{inode} structure.
It creates a per-process table in the file's \texttt{inode->f\_security}.
For each owner domain of a file path, it stores the \emph{domain signature}, \pac{DB}{\texttt{NULL}}{0}, in the table.
If the file is not domain-private, the table is left as \texttt{NULL}.
After, the reference monitor transitions the process into the \emph{protected mode}, where the reference monitor authenticates every path and \fd arguments used in syscalls.


Before authenticating a file object import, the requested file path is first resolved to obtain the underlying \texttt{inode} structure.
This is to avoid confusion when encountering relative paths and symbolic links.
If the table for the current process in \texttt{inode->f\_security} is \texttt{NULL}, the reference monitor recognizes that it is an \emph{ambient} object and can be accessed without further authentication.
If domain signatures are present, the reference monitor authenticates them using the corresponding domain key until a valid signature is found.
The reference monitor verifies that the current user context is in the object's owner domain if any valid signature is found while rejecting the file access otherwise.
Finally, the returned \fd for the requested file is signed with \key{DA} for ambient objects or the domain key for domain-private objects.
% Details of \fd signing are presented in the next subsection.

% \added{

% }


% The reference monitor will intervene and examine the \texttt{inode} structure.


\subsection{Enforcing domain-private file descriptors}
% \subsection{File descriptors as in-process capabilities}
\label{sec:design:fd}
% \begin{figure}[t] \centering \includegraphics[width=0.90\columnwidth]{images/capacity-fd-filter}
% 	\caption{\thename file descriptor authentication logic}
% 	\label{fig:ref}
% \end{figure}
\keypoint{
	\thename's reference monitor intercepts syscalls that generate \fds and create signed \fds before they are sent to the userspace.
	It also intercepts \fd-accepting syscalls (\eg, \texttt{read}) for \fd authentication.
}
Hence, all \fds issued to the process carry \glspl{ac} that are authenticated on use.
Additionally, we use a \emph{secret modifier} only known to the kernel as the \gls{pa} signing modifier, to prevent forging by signing a plain integer with \texttt{pac} instructions in the userspace (\srforge).
\thename's security model also requires that the integer for \fds are not reused (\srreuse), even after they are closed.
We achieve this by changing the hook for the \texttt{close} to reserve the closed \fd in the thread's \fd table, preventing the OS from reusing it.

\hdr{In-process FD capabilities}
\keypoint{
	\thename's \fd access control scheme resembles that of Capsicum's capability-enabled \fds~\cite{capsicum} but is much more lightweight and is enforced at the in-process granularity.
}
\thename's \fd capabilities scheme requires no additional kernel metadata for the permissions of \fds since the information necessary to authorize an \fd use is engraved in the \fd itself.
Moreover, thanks to integrity protection provided by cryptographic authentication, we can securely embed fine-grained access control attributes into the file descriptors themself (\texttt{FD->Attr.} in \autoref{fig:ref-encapsulation}) without extensive bookkeeping.
% and the \fds are integrity-protected by \gls{pa}'s cryptographic authentication.
A combination of \gls{pa}-assisted authentication on the \fd and a per-syscall check based on static policies is sufficient to grant or reject file descriptor access.

Our prototype supports four capabilities, represented as a bitmask of enabled capabilities.
The \texttt{CAP\_WRITE} and \texttt{CAP\_READ} capabilities allow read-related syscalls (\eg, \texttt{read}, \texttt{recvfrom}) and write-related syscalls (\eg, \texttt{write}) to use the \fd.
\texttt{CAP\_SOCKET} enables network-related syscalls, such as \texttt{accept} and \texttt{listen}.
\texttt{CAP\_DELEGATE} allows the \fd to be delegated to other domains.
The reference monitor uses a capability bitmask for each syscall to validate the \fd's attributes.
% The capabilities are revoked from an \fd using simple bit masking operations.
% We also introduce an API to remove capabilities from \fds, 
On the other hand, \texttt{cap\_limit\_fd()} instructs the reference monitor to remove the selected capabilities from an \fd and re-sign it.


\hdr{Signing and authentication of 32-bit \fds}
\autoref{fig:ref-encapsulation} demonstrates the signing strategy of file descriptors.
Since most file descriptors are 32-bit signed integers, we devised a 32-bit signing and authentication scheme.
To sign an \fd, the reference monitor first signs its 64-bit zero-extended value like a normal pointer.
Then, it attaches the PAC to bits[30:23] of the \fd.
The monitor uses bit[22] (\texttt{FD->D}) bit to distinguish between domain-private and ambient \fd, such that domain-private {\fd}s are authenticated using the domain key.
Bit[17:21] stores the bitmask of the previously described \fd capabilities.
bit[31] is unused to prevent the descriptor from being interpreted as an error code (e.g., \texttt{if (fd < 0)}).
\fd authentication on 32-bit \fds happens likewise.

\hdr{Handling special \fd values} During our implementation, we found that special \fd values, such as \texttt{STDIN}, \texttt{STDOUT}, or \texttt{AT\_FDCW}, need to be handled differently since they are often hard-coded into the application as integer values (e.g., 1 and 2). We currently make our reference monitor omit their authentication for compatibility, given that they are always ambient objects.

% \hdr{Considerations for \fd reuse}

% NOTE: 
% \hj{maybe remove this}
% \added{
% 	Note that the limited number of capabilities currently supported by \thename is not a fundamental limitation.
% 	One approach to support more capabilities would be to reduce the number of available bits for file descriptor values to express additional capabilities.
% 	Furthermore, the decision to support 32-bit \fd capabilities was made for compatibility with existing applications.
% 	If a program were to support 64-bit file descriptors, it would be possible to provide more fine-grained capabilities while maintaining a sizable file descriptor space.
% }



% \hdr{Delegation}

% In our example in \autoref{fig:code-example}, \action{FD-Auth} occurs in \texttt{sys\_read} under the monitor's intervention.
% The monitor first checks \texttt{FD->D} and finds that it should be authenticated with the domain key.
% % Then, the reference monitor moves the \gls{pac} value from the 32-bit file descriptor to bits [54:48] of the zero-extended \fd and invokes the authentication instruction on the \fd.
% If authentication succeeds, the requested \gls{syscalsysteml} is then performed on the raw \fd.

% \hdr{Delegating \fds} 
% A domain can also request the reference monitor to create  \fds usable by another domain from an \fd that it owns, using the \texttt{delegate\_fd} API. 
% We further demonstrate the usage of \fd delegation in the \thename-enabled webserver in \cref{sec:evaluation:porting-app}.

% Not terribly important
% \hdr{Special file descriptors and \texttt{stdin}/\texttt{stdout}}
% \thename also supports \fds that are not file-backed like socket, pipe, or shared memory.
% For example, socket \fds are signed and authenticated in \glspl{syscall} (e.g., \texttt{sys\_socket}, \texttt{sys\_accept4}) identically to the case of sensitive files shown above.
% Also, in the current implementation, we heuristically ignore \texttt{stdin},\texttt{stdout}, \texttt{stderr} (\fd numbers 0-3) that are opened by the libc during its initialization, along with other OS-reserved special \fd values.
% This is because those \fd are almost always ambient objects.

% The user can configure whether newly created non-file \fds should be domain-private or ambient; for the former case, our reference monitor always returns domain-private \fds when {\fd}s could create new {\fd}s that are signed by \thename's reference monitor.
% In fact, we demonstrate socket isolation with our real-world application example in \cref{sec:evaluation}.



% Note that \thename's capability model also protects {\fd}s that do not have backing files in the file system, e.g., {\fd}s for sockets and events.
% For instance, socket {\fd}s are protected by signing the return value of \texttt{sys\_socket} \gls{syscall} with the domain key, making it domain-private.
% The reference monitor then simply authenticates the \fd argument of socket-related \glspl{syscall} (e.g., \texttt{sys\_bind}, \texttt{sys\_accept4}) and signs their return \fd returned values.



